% ---------------------------------
\section*{Abstract}

This project ran four AI tools in sequence
to create a knowledge artifact on Privacy
and Data Protection in AI systems, a core
topic from Neeli Prasad's IEEE Distinguished
Lecture on Human-AI Collaboration. The
workflow ran in stages: Perplexity gathered
research on privacy risks and regulatory
frameworks, Claude API synthesized the
findings into a six-page article, QuillBot
rewrote a sample in a different style to
compare clarity, and for the deck ChatGPT
meta-prompted the SlidesAI prompt (GPT-5.5,
Soft Sage template) before SlidesAI.io
rendered a 10-slide presentation from the
curated material. Every prompt,
configuration, and output is documented so
the process can be reproduced and examined.

Claude API produced the clearest technical
writing, and QuillBot showed how a change in
vocabulary shifts readability, while
SlidesAI.io built the slides with no manual
design work. The free tiers did impose
limits, since QuillBot locked its Formal and
Creative modes behind a subscription and
SlidesAI.io offered rigid templates. This
report documents the full methodology,
compares the tools across their different
roles, and reflects on what it means to use
multiple AI systems responsibly in academic
work.
% The approach mirrors Prasad's call for transparency: showing all the steps, not just the final product.
